\section{Verdict Generation}\label{sec:verdicts}
The first matching rule applies: (1) identical sources give no verdict; (2) any confirmed-difference input gives \DIV; (3) any blocking condition gives \INS{} with the lowest-ranked code as primary; (4) otherwise \NDO{} at the highest level whose floors hold (Fig.~\ref{fig:verdict}). Blocking conditions in rank order: import or entry failure; unsupported signature or alignment failure; sandbox failure; unconfirmed or environment-sensitive candidate; quarantined input that may hit a changed unit; exactly one side limited; truncated capture; no usable input or unobservable outcome; an S1 floor failing. The closed reason-code enumeration has 18 codes; four (ambiguous entry, unmodelled change, unsafe program except trace tamper, blind-spot dominance) are never produced by the current engine, so the \INS{} vocabulary is partly aspirational.

Equivalence cannot be emitted because the verdict type asserts a closed value set; this is a naming guarantee. \NDO{} means only that, on the explored deterministic inputs and under $H$, every changed unit ran and no observed-channel difference was found.

\textbf{A defect found and fixed.} An independent review reproduced that an input with no-run outcome (both sides hit a limit) was silently dropped, so \NDO{} was reachable with an incomplete input. The fix (commit 700532b8) makes it a blocker with primary reason \texttt{resource\_limit}. The development run was repeated; three pairs changed from \NDO{} (S2) to \INS{}; all results in this paper are post-fix (Section~\ref{sec:results}). The pre-fix numbers are superseded and not used.
\dotfig{verdict_flowchart}{Verdict flow: three verdicts, first match wins, no equivalence verdict}{fig:verdict}
